\chapter{Supplementary Results}
\label{app:supplementary-results}

\section{Reasoning for the performance of TrP under measurement noise}
\label{sec:app-perf-trp-measurment-noise}

% TODO: Apply tildes to all the inflated parts here

We start by looking at the clean loop. Here, any $y_k$ lies in $\operatorname{range}(U_1)$ by construction.
We also have $\tilde{\Sigma}_0 = U_1 \Sigma_0 U_1\T$. % TODO: Why?
Solving now \eqref{eq:state-correlation}, we obtain $\tilde{\Sigma}_{K_j} = U_1 \Sigma_{K_j} U_1\T$, where $\Sigma_{K_j}$ is the solution
of the uninflated loop.

We now define the projector onto the complement of $\operatorname{range}(U_1)$ as $P_{\perp_1} = I_{d_{\mathrm{in}}} - U_1 U_1\T$.
For any $K_j$, we can therefore seperate it into $K_j = K_jU_1 U_1\T + K_jP_{\perp_1}$. If we now take a look at the updated of this second
component in \cref{alg:mb-pg}, we get $K_{j+1}P_{\perp_1} = K_jP_{\perp_1} - \eta \nabla f(K_j)P_{\perp_1}$. Using the gradient calculation
from \eqref{eq:pg-gradient}, the part on the right turns into 
$\nabla f(K_j)P_{\perp_1} = 2 \left( \left( R + B\T P_{K_j} B \right) K_j - B\T P_{K_j} A \right) \tilde{\Sigma}_{K_j} P_{\perp_1}$, where
$\tilde{\Sigma}_{K_j} P_{\perp_1} = U_1 \Sigma_{K_j} U_1\T (I_{d_{\mathrm{in}}} - U_1 U_1\T) = U_1 \Sigma_{K_j} U_1\T - U_1 \Sigma_{K_j} U_1\T = 0$. % TODO: Break into multilines
Plugging this back into \needcite{ref to gradient calc}, we also have $\nabla f(K_j)P_{\perp_1} = 0$, and therefore
$K_{j+1}P_{\perp_1} = K_jP_{\perp_1} - 0$ holds for all $j$, meaning $K_j P_{\perp_1} = K_0 P_{\perp_1}$, i.e. is never updated
via the policy gradient. And as we construct $\tilde{K}_0 = U_2\T K_0 U_1\T + K_{\perp}$, where 
$\tilde{K}_0 P_{\perp_1} = U_2\T K_0 U_1\T P_{\perp_1} + K_{\perp} P_{\perp_1} = 0 + K_{\perp} P_{\perp_1}$, we know this component
to not be zero.

That the gradient equals zero for this part of the gain also makes sense, it is not associated with any commanded actuation, as
$K_{j} P_{\perp_1}(U_1 x) = 0$. However, under measurement noise, where we have the noise component $n_m$, and $y_k = U_1 x + n_{m_k}$

2. \eqref{eq:pg-gradient} multiplies by $\tilde{\Sigma}_K$ from the right, so with $P_{\perp} = I_{d_{\mathrm{in}}} - U_1 U_1\T$ every gradient
satisfies $\nabla f(\tilde{K}_j) P_{\perp} = 0$. Hence $\tilde{K}_j P_{\perp} = \tilde{K}_0 P_{\perp} = K_{\perp} P_{\perp}$ for all $j$
(\eqref{eq:initial-gain}, as $U_1\T P_{\perp} = 0$): \emph{TrP} never changes the part of the gain that sees only the invisible directions.

3. Without noise this part has no effect on the cost. Under measurement noise ($y_k = U_1 x_k + n^{\mathrm{m}}_k$, Ch. 2 l. 544, noise in all
$d_{\mathrm{in}}$ directions), $\tilde{v}_k$ contains $-K_{\perp} P_{\perp} n^{\mathrm{m}}_k$. It is charged through $\tilde{v}_k\T \tilde{R}\,
\tilde{v}_k$ with $\tilde{R} = I$, and enters the dynamics through $B U_2 K_{\perp} P_{\perp} n^{\mathrm{m}}_k$ ($U_2 K_{\perp} U_1 = 0$ does not
imply $U_2 K_{\perp} P_{\perp} = 0$). With $\sigma_{\mathrm{m}}^{2} I_{d_{\mathrm{in}}}$, the control cost alone grows by at least
$\sigma_{\mathrm{m}}^{2} \lVert K_{\perp} P_{\perp} \rVert_F^2$ per step, a term no iteration of \emph{TrP} reduces.

4. S-PI's proximal step soft-thresholds all singular values, so it shrinks this part as well (state as "can remove", unless shown in general).

5. \emph{IdP}: the same argument on the identified model with $\operatorname{range}(\hat{C})$ in place of $\operatorname{range}(U_1)$.


\section{Additional Plots for the Perception Experiments}
\label{sec:app-perc-plots}

\begin{figure}[htbp]
  \centering
  \generatedfigure{ch6_perception}{ch6_perception_trajectories}
  \caption[Example trajectories of the PPO variants.]{Trajectories of the PPO variants for one sampled episode under one seed.}
  \label{fig:app-perc-trajectories}
\end{figure}


\section{Exit Rates of the Perception Experiments under Noise}
\label{sec:app-perc-noise-terminated}

\Cref{tab:perc-noise-terminated} notes how often the mass leaves the screen under each noise type and
intensity in \cref{subsec:perception-noise-results}.

\begin{table}[htbp]
\centering
\caption{Share of test episodes in the camera environment that end with the mass leaving the view, in percent, per noise channel and intensity. Mean $\pm$ standard deviation over the seeds; \emph{none} is the clean loop.}
\label{tab:perc-noise-terminated}
\begin{tabular}{@{}llrrrrr@{}}
\toprule
Channel & Solver & \emph{none} & \emph{weak} & \emph{standard} & \emph{strong} & \emph{divergent} \\
\midrule
Measurement & FcE & $0$ & $0 \pm 1$ & $3 \pm 2$ & $14 \pm 11$ & $66 \pm 26$ \\
 & OrB & $0$ & $0$ & $3 \pm 7$ & $19 \pm 27$ & $66 \pm 25$ \\
 & ScE & $6 \pm 3$ & $8 \pm 3$ & $10 \pm 5$ & $17 \pm 9$ & $42 \pm 30$ \\
 & PVE & $3 \pm 8$ & $18 \pm 24$ & $43 \pm 41$ & $73 \pm 34$ & $79 \pm 27$ \\
\midrule
Control & FcE & $0$ & $0$ & $0$ & $0$ & $0$ \\
 & OrB & $0$ & $0$ & $0$ & $0$ & $0$ \\
 & ScE & $6 \pm 3$ & $6 \pm 3$ & $7 \pm 4$ & $10 \pm 5$ & $18 \pm 5$ \\
 & PVE & $3 \pm 8$ & $4 \pm 9$ & $4 \pm 9$ & $11 \pm 18$ & $21 \pm 28$ \\
\midrule
Parametric & FcE & $0$ & $8 \pm 3$ & $8 \pm 3$ & $16 \pm 2$ & $42 \pm 5$ \\
 & OrB & $0$ & $8 \pm 3$ & $9 \pm 2$ & $16 \pm 5$ & $39 \pm 5$ \\
 & ScE & $6 \pm 3$ & $52 \pm 6$ & $52 \pm 6$ & $62 \pm 5$ & $74 \pm 3$ \\
 & PVE & $3 \pm 8$ & $36 \pm 22$ & $37 \pm 23$ & $48 \pm 20$ & $62 \pm 13$ \\
\midrule
Disturbance & FcE & $0$ & $0$ & $0$ & $0 \pm 1$ & $89 \pm 3$ \\
 & OrB & $0$ & $0$ & $0$ & $5 \pm 3$ & $85 \pm 5$ \\
 & ScE & $6 \pm 3$ & $9 \pm 4$ & $10 \pm 5$ & $18 \pm 5$ & $46 \pm 7$ \\
 & PVE & $3 \pm 8$ & $3 \pm 8$ & $4 \pm 8$ & $8 \pm 10$ & $36 \pm 15$ \\
\bottomrule
\end{tabular}
\end{table}

